\section{Introduction}
\label{sec:introduction}
\subsection{Problem and motivation}
Rain, fog, snow and smoke obscure objects and reduce contrast. Restoration
can remove corruption, but it can also erase boundaries or lose detail during
resizing. Higher PSNR therefore need not improve detection.

Aegis Perception tests that distinction through paired-image benchmarks,
fixed-detector comparisons and execution on the original 4\,GB Jetson Nano.
The project follows one selected checkpoint from reconstruction to downstream
accuracy and processing cost (Fig.~\ref{fig:project-workflow}).

\subsection{Research questions and scope}
\begin{enumerate}
\item Which of six restoration recipes leads a shared four-weather benchmark?
\item Do model or loss changes outperform equal-budget continued training?
\item Does restoration help a fixed detector after controlling input geometry?
\item What accuracy, throughput and energy does the Nano cascade deliver?
\end{enumerate}
These questions produce six completed work packages (Fig.~\ref{fig:timeline}).
All models use reconstruction objectives; detector training and detector-loss
optimization lie outside the completed study. Different pretrained starting
points mean that the benchmark compares recorded recipes, not architecture alone.

\subsection{Main contribution}
The project connects an audited data split, six trained models, eight matched
continuation arms and measured detection/device results. The rejected
continuation changes and the Nano accuracy loss establish useful limits on
this pipeline. Supplementary Material preserves detailed controls, validation
trajectories and the evidence map.
